% !TEX root = main.tex
\section{Extend the Field of Values for the Collatz Sequence}\label{sec:extend-the-field-of-values-for-the-collatz-sequence}

A better thought is extending our view from $\mathbb{N}$ to $\mathbb{Q}$.
Keeping $\alpha, \beta, \gamma \in \mathbb{N}$, instead of restricting the initial number $a \in \mathbb{N}$, we can let $a$ be any number of $\mathbb{Q}-\beta\mathbb{Q}$.
In this situation, the \cref{theorem:inductive-formula} still holds, so do those corollaries of the theorem.
More than that, we can get a Collatz sequence by given the reduces.

\begin{theorem}~\label{theorem:any-cycle}
    Given $\{v_{i}\}_{i=1}^{n}$ that $n < \infty$, if $\forall i>j\ge1, \beta \nmid \mathfrak{A}_i^{j}$, then there exists only one Collatz sequence in $\mathbb{Q}$ is a cycle that $\val_{\beta}(a;k) = v_{k\fmod n}$.
\end{theorem}
\begin{proof}%[Proof of \cref{theorem:any-cycle}]
    Let's just using the notation $\mathcal{A}_k$ and $\mathfrak{A}_k$ as if $\val_{\beta}(a;k) = v_{k\fmod n}$.\newline
    Then $\forall k = 1, 2, \cdots, n$, we have $\mathcal{A}_{k + n}^{k} = \mathcal{A}_n$.
    Let
    \[ a_k = \frac{\gamma\cdot\mathfrak{A}_{n+k}^{k}}{\beta^{\mathcal{A}_n} - \alpha^n} \in \mathbb{Q}. \]
    Obviously, we have $\beta \nmid a_{k}$.
    And we have
    \begin{align*}
        \alpha\cdot a_{k} + \gamma &= \frac{\gamma\cdot\alpha\cdot\mathfrak{A}_{n+k}^{k}}{\beta^{\mathcal{A}_n} - \alpha^n} + \frac{\gamma\cdot(\beta^{\mathcal{A}_n} - \alpha^n)}{\beta^{\mathcal{A}_n} - \alpha^n} \\
        &= \frac{\gamma\cdot(\alpha\cdot\mathfrak{A}_{n+k}^{k}+\beta^{\mathcal{A}_{n+k}^{k}}) - \gamma\cdot\alpha^n}{\beta^{\mathcal{A}_n} - \alpha^n} \\
        &= \frac{\gamma\cdot(\mathfrak{A}_{n+k+1}^{k} - \alpha^{n})}{\beta^{\mathcal{A}_n} - \alpha^n} \\
        &= \frac{\gamma\cdot\beta^{\mathcal{A}_{k+1}^{k}}\cdot\mathfrak{A}_{n+k+1}^{k+1}}{\beta^{\mathcal{A}_n} - \alpha^n} \\
        &= \frac{\gamma\cdot\mathfrak{A}_{n+k+1}^{k+1}}{\beta^{\mathcal{A}_n} - \alpha^n}\cdot \beta^{\mathcal{A}_{k+1}^{k}} \\
        &= a_{k+1}\cdot \beta^{\Delta_{k+1}}.
    \end{align*}
    So we just let $a = a_{n}$, then the Collatz sequence of $a$ is what we want.
    And it is unique.
\end{proof}
In fact, we can see that the rotation of $v_{k}$ corresponds with the rotation of $a_k$.
That implies some symmetry of the Collatz sequence.
But not all Collatz reduce is a pure cycle.
In fact, most of  what may eventually enter a cycle have a chaotic beginning.
Even there may exist some Collatz sequence will not fall into any cycle.
For these two kinds of Collatz reduces, we have
\begin{theorem}~\label{theorem:any-presequence}
    Given $\{v_{i}\}_{i=1}^{n}$ that $n < \infty$, and any number $b$.
    If $\forall i>j\ge1, \beta \nmid \mathfrak{A}_i^{j}$, then there exists only one Collatz sequence in $\mathbb{Q}[b]$ that $\val_{\beta}(a;i) = v_{i}$ and $a_n = b$.
\end{theorem}
\begin{proof}
    Also using the notation $\mathcal{A}_k$, $\mathfrak{A}_k$ and $\mathfrak{a}_k$ as if $\val_{\beta}(a;i) = v_{i}$.
    Let
    \[ a_k = \frac{b \cdot \beta^{\mathcal{A}_n^{k}} - \gamma\cdot\mathfrak{A}_n^{k}}{\alpha^{n-k}} = b \cdot \frac{\beta^{\mathcal{A}_n^{k}}}{\alpha^{n-k}}  - \frac{\gamma}{\alpha}\cdot\alpha^{k}\cdot\mathfrak{a}_n^{k} \in \mathbb{Q} \]
    where $k = 0, 1, 2, \cdots, n$.
    Similarly, since $\beta \nmid \mathfrak{A}_n^k$, then $\beta \nmid \mathfrak{a}_n^{k}$, then $\beta \nmid a_k$.
    And we have
    \begin{align*}
        \alpha \cdot a_k + \gamma &= b \cdot \frac{\beta^{\mathcal{A}_n^{k}}}{\alpha^{n-k-1}}  - \gamma\cdot\alpha^{k}\cdot\mathfrak{a}_n^{k} + \gamma \\
        % &= b \cdot \frac{\beta^{\mathcal{A}_n^{k}}}{\alpha^{n-k-1}}  - \gamma\cdot\alpha^{k}\cdot\left(\mathfrak{a}_n^{k} - \frac{1}{\alpha^{k}}\right) \\
        &= b \cdot \frac{\beta^{\mathcal{A}_n^{k}}}{\alpha^{n-k-1}}  - \gamma\cdot\alpha^{k}\cdot\left(\mathfrak{a}_n^{k} - \mathfrak{a}_{k+1}^{k}\right) \\
        &= \left(b \cdot \frac{\beta^{\mathcal{A}_n^{k+1}}}{\alpha^{n-(k+1)}}  - \frac{\gamma}{\alpha}\cdot\alpha^{k+1}\cdot\mathfrak{a}_{n}^{k+1}\right) \cdot \beta^{\mathcal{A}_{k+1}^{k}} \\
        &= a_{k+1} \cdot \beta^{\Delta_{k+1}}
    \end{align*}
    to make sure it is a Collatz sequence.
    Finally, we have $a_n = b$.
    Also, it is unique.
\end{proof}
\begin{corollary}
    Given $a = \frac{   p}{q}\in \mathbb{Q}$. If $\exists s\in \mathbb{N}$, such that $q = \alpha^{s}$, then $a_{s} \in \mathbb{N}$ by natual reduce sequence.
\end{corollary}

\cref{theorem:any-cycle} and \cref{theorem:any-presequence} present that we can come up with a lot of Collatz sequence that will fall into cycles by given its reduce sequence instead of startswith an integer.
However, not all given reduce sequence can derive a Collatz sequence because there may be some $\beta\mid \mathfrak{A}_{n+k}^{k}$.
At that time, we have $\beta \left\vert \frac{\gamma\cdot\mathfrak{A}_{n+k}^{k}}{\beta^{\mathcal{A}_n} - \alpha^n} \right.$.
That's why we need $\beta \nmid \mathfrak{A}_i^j$.
Moreover, from the two proofs above, we can also draw another interesting corollary:
\begin{theorem}
    If Collatz sequence of $a$ will fall into a cycle, then $a \in \mathbb{Q}$.
\end{theorem}
However, the converse of this conclusion does not hold.
For example, $a=3$ under $\alpha = 11$, $\beta=2$ and $\gamma=1$ may never fall into a cycle.
For these won't fall in to cycles, there is also a conjecture:
\begin{conjecture}
    Given $\{v_{i}\}_{i=1}^{\infty}$ has no cycle, then there is a Collatz sequence that $\val_{\beta}(a;i) = v_{i}$, it has a subsequence $\{a_{k_i}\}_{i=0}^{\infty}$ converges to an $\beta$-adic irrational number or $0$.
\end{conjecture}

Anyway, we now finally extended the problem from integers to the rational numbers, and at the same time, the exploration of Collatz sequences has turned to the exploration of Collatz reduces.
This implies that the Collatz reduces is the more fundamental aspect of this process.
